
\documentclass{article}
\usepackage{colortbl}
\usepackage{makecell}
\usepackage{multirow}
\usepackage{supertabular}

\begin{document}

\newcounter{utterance}

\centering \large Interaction Transcript for game `cladder', experiment `full\_v1.5\_default', episode 4979 with qwen3.5{-}2b.
\vspace{24pt}

{ \footnotesize  \setcounter{utterance}{1}
\setlength{\tabcolsep}{0pt}
\begin{supertabular}{c@{$\;$}|p{.15\linewidth}@{}p{.15\linewidth}p{.15\linewidth}p{.15\linewidth}p{.15\linewidth}p{.15\linewidth}}
    \# & \multicolumn{2}{c}{Player} && \multicolumn{2}{c}{Game Master} \\
    \hline

    \theutterance \stepcounter{utterance}  
    & & & \multicolumn{4}{p{0.6\linewidth}}{
        \cellcolor[rgb]{0.9,0.9,0.9}{
            \makecell[{{p{\linewidth}}}]{
                \texttt{\tiny{[P1$\langle$GM]}}
                \texttt{You are an expert in causal inference. The following question is not a typical commonsense query, but rather a meticulously designed question created by a professor specializing in causal inference, intended to assess the students' mastery of the course content.} \\
\texttt{Imagine a self{-}contained, hypothetical world with only the following conditions, and without any unmentioned factors or causal relationships: Vaccination status has a direct effect on vaccination reaction and getting smallpox. Getting smallpox has a direct effect on black hair. Vaccination reaction has a direct effect on black hair. For unvaccinated individuals, the probability of black hair is 27\%. For vaccinated individuals, the probability of black hair is 37\%. Will vaccination decrease the chance of black hair?} \\
\texttt{Start your answer with 'Yes' or 'No', followed by additional reasoning or evidence to support your explanation.} \\
            }
        }
    }
     \\ \\

    \theutterance \stepcounter{utterance}  
    & \multicolumn{4}{p{0.6\linewidth}}{
        \cellcolor[rgb]{0.561,0.737,0.561}{
            \makecell[{{p{\linewidth}}}]{
                \texttt{\tiny{[P1$\rangle$GM]}}
                \texttt{Yes} \\
\\ 
\texttt{**Reasoning:**} \\
\texttt{In this hypothetical world, the causal chain is defined as follows:} \\
\texttt{1.  **Vaccination** \$\textbackslash{}rightarrow\$ **Vaccination Reaction** (Direct Effect).} \\
\texttt{2.  **Vaccination Reaction** \$\textbackslash{}rightarrow\$ **Black Hair** (Direct Effect).} \\
\texttt{3.  **Getting Smallpox** \$\textbackslash{}rightarrow\$ **Black Hair** (Direct Effect).} \\
\\ 
\texttt{The question asks whether vaccination decreases the chance of black hair. To determine this, we must look at the net effect of the intervention "Vaccinate" on the outcome "Black Hair," holding all other variables constant.} \\
\\ 
\texttt{*   **Path 1 (Vaccination \$\textbackslash{}rightarrow\$ Reaction \$\textbackslash{}rightarrow\$ Black Hair):** Vaccination increases the probability of black hair by 10 percentage points (from 27\% to 37\%).} \\
\texttt{*   **Path 2 (Smallpox \$\textbackslash{}rightarrow\$ Black Hair):** This path is independent of vaccination status in this model.} \\
\\ 
\texttt{Since the only causal path linking vaccination to black hair is through the vaccination reaction, and this path results in an increase in black hair probability, the net effect of vaccination on black hair is an increase. Therefore, vaccination does **not** decrease the chance of black hair; it increases it.} \\
            }
        }
    }
    & & \\ \\

    \theutterance \stepcounter{utterance}  
    & & & \multicolumn{2}{p{0.3\linewidth}}{
        \cellcolor[rgb]{0.9,0.9,0.9}{
            \makecell[{{p{\linewidth}}}]{
                \texttt{\tiny{[GM$|$GM]}}
                \texttt{Yes} \\
            }
        }
    }
    & & \\ \\

    \theutterance \stepcounter{utterance}  
    & & & \multicolumn{2}{p{0.3\linewidth}}{
        \cellcolor[rgb]{0.9,0.9,0.9}{
            \makecell[{{p{\linewidth}}}]{
                \texttt{\tiny{[GM$|$GM]}}
                \texttt{no} \\
            }
        }
    }
    & & \\ \\

    \theutterance \stepcounter{utterance}  
    & & & \multicolumn{2}{p{0.3\linewidth}}{
        \cellcolor[rgb]{0.9,0.9,0.9}{
            \makecell[{{p{\linewidth}}}]{
                \texttt{\tiny{[GM$|$GM]}}
                \texttt{game\_result = LOSE} \\
            }
        }
    }
    & & \\ \\

\end{supertabular}
}

\end{document}
